\chapter*{前\hspace{1em}言}
\markboth{前\hspace{1em}言}{前\hspace{1em}言}
\phantomsection\label{ch:preface}

本书按编写、检查、运行和交付程序的顺序介绍 ChaLi。先通过小程序建立类型与控制流的基本概念，再学习值、引用、资源、泛型和接口，最后组合成数据库应用。正文描述当前 1.0.11 已实现的行为；版本号以仓库 \href{https://github.com/rabbitc999/ChaLi/blob/master/VERSION}{VERSION} 为准，未完成计划不作为语法介绍。

\section*{阅读与运行}

没有编程经验的读者可从第 1 章开始。已有经验的读者可借助目录查找专题；第 6 至第 8 章的值、引用和资源规则建议一起阅读。每章末尾提供习题与参考解答。中文番外及技术对照放在章末，英文版只包含技术正文和附录。

完整程序带 main 和预期输出。交互例程另列固定测试输入，终端的输入回显不计入程序输出。短片段用于展示声明或表达式，需要放入相应上下文。错误示例用于说明被拒绝的条件，不应与正常例程合并运行。

命令中的 chali 表示已加入 PATH 的编译器。Windows 源码构建产物带版本、LLVM 变体和平台后缀，执行时换成实际路径；Linux 当前目录的编译器写作 ./chali。安装与平台范围见\href{https://github.com/rabbitc999/ChaLi/blob/master/doc/39-installation-packages.md}{安装说明}。应用使用的数据库与动态库需单独准备。

可以在线阅读 Markdown，或使用 \href{https://github.com/rabbitc999/ChaLi/blob/master/doc/book/latex/book.pdf}{PDF}。\href{https://github.com/rabbitc999/ChaLi/blob/master/doc/book/en/README.md}{英文版} 与中文技术正文对应，不包含番外，也不生成英文 LaTeX/PDF。

\section*{章节地图}

\begin{center}
\begin{tabular}{lll}
\toprule
章 & 标题 & 内容 \\
\midrule
\hyperref[ch:01]{01} & 第一个程序 & 编译工具、检查、运行与导出 \\
\hyperref[ch:02]{02} & 猜数字 & 固定输入的交互循环 \\
\hyperref[ch:03]{03} & 变量与类型 & 声明、数值类型与显式转换 \\
\hyperref[ch:04]{04} & 函数 & 参数、返回值和调用 \\
\hyperref[ch:05]{05} & 控制流 & 分支、循环与控制表达式 \\
\hyperref[ch:06]{06} & 值语义 & 复制、借用、移动和可观察行为 \\
\hyperref[ch:07]{07} & 引用 & 引用有效期及可变访问 \\
\hyperref[ch:08]{08} & 资源与作用域 & 作用域退出、延迟清理和资源 \\
\hyperref[ch:09]{09} & 结构体 & 字段、成员方法和具名构造入口 \\
\hyperref[ch:10]{10} & 枚举与模式匹配 & 状态、载荷与穷尽匹配 \\
\hyperref[ch:11]{11} & 容器 & 数组、List、Map、Set 与容量 \\
\hyperref[ch:12]{12} & 字符串与文本 & UTF-8、格式化及安全输出 \\
\hyperref[ch:13]{13} & 接口与运行期多态 & 行为抽象、继承和具体类型匹配 \\
\hyperref[ch:14]{14} & 泛型 & 类型参数、常量参数和约束 \\
\hyperref[ch:15]{15} & 函数值与闭包 & 环境捕获、返回闭包与调用边界 \\
\hyperref[ch:16]{16} & 错误处理 & Option、Result、接口载荷与传播 \\
\hyperref[ch:17]{17} & 迭代器 & 消费状态、组合与结构修改 \\
\hyperref[ch:18]{18} & 模块与工程 & 模块归属、工程目录和配置 \\
\hyperref[ch:19]{19} & 预处理与宏 & 卫生、参数类别与结果类型推导 \\
\hyperref[ch:20]{20} & 注解与反射 & 字段元数据与序列化边界 \\
\hyperref[ch:21]{21} & 与 C 互操作 & 原型、封送、资源与同步回调 \\
\hyperref[ch:22]{22} & 编译器指令 & 派生、运算和借用来源 \\
\hyperref[ch:23]{23} & 不安全代码 & 地址、初始化状态与原始存储 \\
\hyperref[ch:24]{24} & AOT 导出 & LLVM、目标配置和部署依赖 \\
\hyperref[ch:25]{25} & 学生信息管理系统 & 登录、权限与教务数据的综合应用 \\
\bottomrule
\end{tabular}
\end{center}

\section*{附录与参考}

\begin{center}
\begin{tabular}{ll}
\toprule
入口 & 用途 \\
\midrule
\hyperref[ch:appA]{附录 A} & 按诊断阶段定位常见错误 \\
\hyperref[ch:appB]{附录 B} & 按任务查找模块与 API \\
\href{https://github.com/rabbitc999/ChaLi/blob/master/doc/04-language-reference.md}{语言参考} & 声明与表达式的详细写法 \\
\href{https://github.com/rabbitc999/ChaLi/blob/master/doc/library/README.md}{标准库参考} & 函数签名、返回值与限制 \\
\href{https://github.com/rabbitc999/ChaLi/blob/master/doc/17-capabilities.md}{能力说明} & 当前实现范围 \\
\href{https://github.com/rabbitc999/ChaLi/tree/master/example/samples}{运行示例} & 仓库内的独立程序 \\
\bottomrule
\end{tabular}
\end{center}

教程解释常见组合，参考文档列出接口细节。遇到错误时先定位是解析、类型与借用、原生链接还是应用运行阶段，再核对对应条件和诊断；不要只通过关闭检查绕过问题。
